% !TEX root = ../main.tex

%----------------------------------------------------------------------------------------
% CHAPTER 6 - ESTUDIS I DECISIONS
%----------------------------------------------------------------------------------------

\chapter{Estudis i decisions}
\label{Ch:Studies}

\section{Estudi de tecnologies}

Abans de decidir les tecnologies finals, es va realitzar un estudi exhaustiu de les alternatives disponibles, valorant maduresa, suport comunitari, adaptabilitat al problema i compatibilitat amb llicències obertes.

\subsection{Resum d'avaluacions}

\begin{table}[h]
\centering
\caption{Comparativa de tecnologies avaluades}
\begin{tabular}{|l|l|l|p{4cm}|}
\hline
\textbf{Categoria} & \textbf{Seleccionada} & \textbf{Descartades} & \textbf{Motiu de la selecció} \\
\hline
Framework frontend & React 18 & Vue 3, Angular, Svelte & Arquitectura de components, ecosistema abundant, característiques concurrents \\
\hline
Bundler & Vite 5 & Webpack 5, CRA & Velocitat de compilació, configuració mínima, ESM nadiu \\
\hline
Visualització & Recharts 2 & Chart.js, D3.js & Integració nativa amb React, configuració via props, precisió de D3 \\
\hline
Mapes & Leaflet 1 & Mapbox GL, OpenLayers & Lleugeresa (38KB), simplicitat, llicència BSD \\
\hline
ETL & Node.js 18 & Python 3, Deno & I/O asíncrona, ecosistema npm, coherència amb frontend \\
\hline
CI/CD & GitHub Actions & GitLab CI, Jenkins & Integració nativa, allotjament gratuït, YAML declaratiu \\
\hline
\end{tabular}
\end{table}

\subsection{Detalls de les avaluacions}

\subsubsection{Frameworks frontend}
\begin{itemize}
    \item \textbf{React 18}: Triat per la seva arquitectura basada en components, característiques concurrents (automatic batching, startTransition), i ecosistema abundant amb biblioteques com Recharts i React Leaflet.
    \item \textbf{Vue 3}: Considerada per la seva sintaxi intuïtiva, però el seu ecosistema és més petit i la integració amb Recharts no és tan directa.
    \item \textbf{Angular}: Descartada per la seva complexitat i sobrecàrrega (sistema de dependències, compilador AOT) excessiva per a les necessitats del projecte.
    \item \textbf{Svelte}: Admirada per la seva innovació (compilació a JS pur), però el seu ecosistema reduït i la menor familiaritat van fer decantar per React.
\end{itemize}

\subsubsection{Eines de construcció}
\begin{itemize}
    \item \textbf{Vite}: Seleccionat per la velocitat de compilació (ESM nadiu), configuració mínima, i rendiment de producció comparable a webpack.
    \item \textbf{Webpack 5}: Poderós però amb complexitat de configuració i temps de compilació més llarg.
    \item \textbf{CRA}: Descartada per dependre de webpack i mancar de flexibilitat.
\end{itemize}

\subsubsection{Biblioteques de visualització}
\begin{itemize}
    \item \textbf{Recharts}: Construït sobre React, configurable via props, utilitza D3.js per càlculs d'escala.
    \item \textbf{Chart.js}: Popular però requereix envoltors addicionals per a React i menys flexibilitat que SVG.
    \item \textbf{D3.js direct}: Màxim control però gestió manual complexa del DOM i integració amb React.
\end{itemize}

\subsubsection{Biblioteques de mapes}
\begin{itemize}
    \item \textbf{Leaflet}: Seleccionat per lleugeresa (38KB), simplicitat, i llicència BSD completa.
    \item \textbf{Mapbox GL JS}: Capacitats avançades però pes significativament major (>200KB) i requeriment de token.
    \item \textbf{OpenLayers}: Poderós però excessivament complex per a necessitats bàsiques de visualització de punts.
\end{itemize}

\subsubsection{Tecnologies ETL}
\begin{itemize}
    \item \textbf{Node.js 18}: Seleccionat per I/O asíncrona, ecosistema npm madur, i coherència amb frontend JavaScript.
    \item \textbf{Python 3}: Alternativa popular amb pandas, però integració menys directa amb frontend React.
    \item \textbf{Deno}: Innovador però immadur comparat amb Node.js i amb ecosistema més reduït.
\end{itemize}

\subsubsection{Eines CI/CD}
\begin{itemize}
    \item \textbf{GitHub Actions}: Seleccionat per integració nativa, sintaxi YAML, i allotjament gratuït (2.000 min/mes).
    \item \textbf{GitLab CI}: Potent però requereix migrar el codi a GitLab.
    \item \textbf{Jenkins}: Extremadament flexible però amb sobrecàrrega operativa excessiva.
    \item \textbf{Cron jobs}: Descartat per la necessitat d'infraestructura addicional.
\end{itemize}

%----------------------------------------------------------------------------------------

\section{Decisions de disseny}

\subsection{Arquitectura basada en components}
Es va decidir organitzar el frontend amb components React independents per a reutilització, mantenibilitat i proves aïllades.

\subsection{Gestió d'estat amb hooks de React}
Sobre alternatives com Redux o Context API, els hooks de React (useState, useEffect) ofereixen simplicitat, adaptació natural a components funcionals, i complexitat reduïda sense necessitat de dependències externes.

\subsection{Estilització SCSS amb BEM}
La metodologia BEM amb variables centralitzades a \texttt{\_variables.scss} millora la mantenibilitat i evita col·lissions de classes.

\subsection{Disseny responsive mobile-first}
CSS Grid per a layouts principals i Flexbox per a alineació interna, assegurant experiència òptima en tots els dispositius.

\subsection{Patró de pipeline modular (ETL)}
L'ETL es divideix en tres fases independents (Extract, Transform, Load) que comuniquen a través d'interfícies ben definides, facilitant proves i mantenibilitat.

\subsection{Independència ETL/frontend}
La comunicació exclusiva a través de fitxers JSON al repositori Git proporciona robustesa, flexibilitat de desplegament, i escalabilitat diferencial.

\subsection{Configuració externalitzada}
Variables d'entorn amb \texttt{.env} per a portabilitat entre entorns i millorar la seguretat.

\subsection{Maneig d'errors centralitzat}
Sistema de logging amb nivells diferents (error, warn, info, debug) i gestió granular d'errors que no atura tot el pipeline per fallades individuals.

%----------------------------------------------------------------------------------------
